% Abstract final — SCM paper (MNRAS submission)
% Numbers locked: p=1.19e-10, 53/53, delta_RMSE_out=-7.3
% Source: results/scm_oos/oos_seed_summary.csv (seed 42)

\begin{abstract}
We present the SCM (Shear-Corrected MOND) framework, in which a local
friction term $F_3$---derived from the radial gradient of the
baryon-deficit field---extends the standard baryonic scaling relations
for galaxy rotation curves.

Using the SPARC sample, we evaluate the predictive performance of the
SCM model relative to a baseline MOND-like relation.  We perform
out-of-sample validation using repeated 70/30 galaxy-level splits,
ensuring strict separation between training and test sets.

We find that the SCM model provides a systematic improvement in
predictive accuracy across all test galaxies considered ($53/53$),
with a median out-of-sample improvement of
$\Delta\mathrm{RMSE}_{\rm out} \approx -7.3$ (in proxy units).
A one-sided Wilcoxon test confirms the statistical significance of
this effect ($p = 1.19\times10^{-10}$).

We interpret this result as evidence that the $F_3$ term captures
additional dynamical information not encoded in the baseline relation.
However, we emphasise that the present analysis is phenomenological
and does not assume a specific microphysical origin.

These results demonstrate that local radial gradients carry
predictive information for galaxy dynamics and motivate further
tests on extended datasets and hierarchical validation schemes.
\end{abstract}
